\subsection{Function Notation} 
When a relation is a function, we can think of the first entry of each pair as an ``input'' and the second entry as an ``output,'' and think of the function as a ``process'' that the variable undergoes.
\begin{definition}[Function Notation]
If the relation $R$ is a function, we can write $$R(x)$$ to indicate the value of $y$ such that $(x,y)\in R$.
\end{definition}
\begin{example}
Confirm that $R=\Set{(-7,-8),(0,3),(1,-7),(5,-8)}$ is a function. Then, evaluate each expression.
\begin{itemize}
\item \makeblankfill[\linewidth]{$R$ is a function because}
\item \makeblankfill[\linewidth]{$R(1)=\makeblanksmall$ because}
\item \makeblankfill[\linewidth]{$R(-7)=\makeblanksmall$ because}
\item \makeblankfill[\linewidth]{$R(4)$ \makeblank[1in] because}
\end{itemize}
\end{example}
\begin{example}
The relation $g$ is graphed below. Confirm that $g$ is a function. Then, evaluate each expression.
\begin{center}
\begin{tikzpicture}
\useasboundingbox (-5,-5) rectangle (5,5);
\setgraph{5}{5}{5}{5}
\GraphOnes{}{}
\draw[graph,<-] (-5,-1.5) -- (2,2);
\draw[graph,->] (2,0) -- (5,-3);
\foreach \y in {2,0} \draw[dotempty] (2,\y) circle (\dotsize);
\end{tikzpicture}
\end{center}
\begin{itemize}
\item \makeblankfill[\linewidth]{$g$ is a function because}
\item \makeblankfill[\linewidth]{$g(3)=\makeblanksmall$ because}
\item \makeblankfill[\linewidth]{$g(-4)=\makeblanksmall$ because}
\item \makeblankfill[\linewidth]{$g(2)$ \makeblank[1in] because}
\end{itemize}
\end{example}
